%% kapa: exam=1 round=30 year=2019 period=1 form=A subject=civil_law subject_ko=민법 seq=31 number=31
%% source: Q-Net 감정평가사 기출문제 (cst00309) artlSeq=5205205 file=extracted/1008/5205205/2019년 제30회 감정평가사 1차시험 문제지/2019년 제30회 감정평가사 1차시험 문제지/2019년도 제30회 감정평가사 1차 1교시 A형.hwp
%% conversion: kordoc 4.24.51 -> tools/gen1.py (auto); review=proofread-sonnet; flags=none
\begin{question}{31}
취득시효에 관한 설명으로 옳지 않은 것은? (다툼이 있으면 판례에 따름)
\begin{choices}
  \item 비법인사단은 시효취득의 주체가 될 수 없다.
  \item 부동산의 소유자로 등기한 자가 10년간 소유의 의사로 평온, 공연하게 선의이며 과실없이 그 부동산을 점유한 때에는 소유권을 취득한다.
  \item 10년간 소유의 의사로 평온, 공연하게 동산을 점유한 자는 그 소유권을 취득한다.
  \item 부동산 점유취득시효가 완성되면 점유자는 원칙적으로 시효기간 만료 당시의 토지소유자에 대하여 소유권이전등기청구권을 취득하는데, 이는 채권적 청구권이다.
  \item 공유지분의 일부에 대해서도 시효취득이 가능하지만, 집합건물의 공용부분은 점유취득시효에 의한 소유권취득의 대상이 될 수 없다.
\end{choices}
\end{question}
\answer{1}
